\chapter{Introdução ao Diagnóstico Energético}

Este relatório técnico apresenta o diagnóstico energético da unidade consumidora <<nomeUc>>, vinculada à Prefeitura Municipal de Salvador. O objetivo principal deste trabalho é identificar oportunidades de racionalização e uso eficiente da energia elétrica, contribuindo para a redução de custos operacionais e para a sustentabilidade ambiental da instituição.

A metodologia adotada inclui a realização de visita técnica in loco, análise do histórico de consumo da unidade e levantamento detalhado dos sistemas e equipamentos elétricos instalados. As ações propostas são fundamentadas nas diretrizes estabelecidas pelo \ac{PEE} da \ac{ANEEL}, conforme regulamentado pela Resolução Normativa nº 920/2021~\cite{aneel920}, e seguem as orientações do \ac{PROPEE}.

A seguir, são apresentados os dados cadastrais da unidade consumidora analisada, os quais servem de base para o planejamento das intervenções e ações corretivas.

\vspace{1cm}



\begin{table}[h]
\centering
\begin{tabular}{|>{\columncolor[HTML]{1a3a5c}}p{5cm}|p{8cm}|}
\hline
\textcolor{white}{\textbf{Nome da Unidade}} & <<nomeUc>> \\
\hline
\textcolor{white}{\textbf{Código da UC}} & <<codigoUc>> \\
\hline
\textcolor{white}{\textbf{Tipo de Atividade}} & <<tipoAtividade>> \\
\hline
\textcolor{white}{\textbf{Endereço Completo}} & <<endereco>> \\
\hline
\textcolor{white}{\textbf{Responsável / Cargo}} & <<responsavel>> \\
\hline
\textcolor{white}{\textbf{Telefone para Contato}} & <<telefone>> \\
\hline
\textcolor{white}{\textbf{E-mail Institucional}} & <<email>> \\
\hline
\end{tabular}
\caption{Dados cadastrais da unidade consumidora avaliada}
\label{tab:dados_uc}
\end{table}







%\begin{table}[h]
    %\centering
    %\includegraphics[width=0.95\textwidth]{Figuras/Tabela 1.png} % substitua pelo nome correto da imagem
    %\caption{Dados cadastrais da unidade consumidora avaliada}
   % \label{tab:dados_uc}
%\end{table}





